% This is a file that contains the tracking of the activities related to the project during the third week of March 2025.

\documentclass[12pt]{article}

\usepackage[a4paper, margin=1in]{geometry}
\usepackage{graphicx}
\usepackage[backend=biber,style=numeric]{biblatex}
\usepackage[most]{tcolorbox}

\addbibresource{references.bib}
\newcommand{\apj}{The Astrophysical Journal}

\begin{document}

\author{Joan Ronquillo}

\title{Project Progress - Week 3 of March 2025}
\maketitle

This is a file that contains the tracking of the activities related to the OFC project during the third week of March 2025.
The activities are divided into groups and a summary of the progress is provided.

\section{Initial Status of the Project}
The project has made significant progress up until the third week of March 2025. The key achievements include:
\begin{itemize}
    \item The management of Python modules and compilation of LaTeX projects has been learned.
    \item A first approach to C++ modular development of the functionalities has been initiated.
    \item The reading of Draine's paper on radiative corrections in polarizability \cite{Draine1988ApJ...333..848D} has been initiated.
    \item The difficulties in implementing C++ modules have been discussed with Raúl. And intuition on how to proceed has been provided.
    \item The need to use the Discrete Dipole Approximation (DDA) for the calculation of optical forces for particles of size comparable to the wavelength has been discussed.
\end{itemize}
So far, the project is developing in some areas:

\begin{itemize}
    \item The development of the algorithm to obtain the polarizability of the particles.
    \item The modular development, including makefile documentation and Raul's cpp repository and notes analysis.
    \item The development of a calculation method for Optical Forces given a polarizability, a field, and a field gradient.
\end{itemize}

\section{Potential Tasks for the Week}
The following tasks are proposed for the week of March 10th to March 15th, 2025:
\begin{itemize}
    \item Exploration of CMakefile, include directory, and modular development (Notion) 
    \item Develop the algorithm to obtain the polarizability of the particles.
    \item Implement the calculation method for Optical Forces given a polarizability, a field, and a field gradient.
    \item Documentation of quasiestatic polarizability theory.
    \item Continue reading Draine's paper on radiative corrections in polarizability \cite{Draine1988ApJ...333..848D}.
    \item Analysys of Raúl's \path{cpp-template-project} repository.
    \item Check of Raúl's class notes.
    \item Discuss the progress and difficulties with Raúl.
\end{itemize}

\section{Week Progress}

\subsection{Monday, March 17th, 2025}
The tasks for the week have been defined and the progress
of the project has been analyzed. CMakeFiles and include
directories have been explored, creating an initial structure
(it doesn't compile yet) and notion pages for registering
information about the template project \path{cpp-template-project}.

\subsection{Tuesday, March 18th, 2025}
The compilation issues with the CMakeFiles have been solved
and the environment dependencies had been installed. The
first test for optical forces calculation function has been
implemented and tested. The function is able to calculate
the optical forces given a medium permittivity, a set of
dipolar momenta and a set of electric field gradients.
A library for dielectric coefficients has been created.
In it, a class for permittivities has been defined. A method
for gold particles testing has been implemented, but yet not
tested correctly.

\section{Next Steps}
The next steps for the project are:
\begin{itemize}
    \item Check the wavelenght dependence of gold's permittivity.
    \item Explore compilator debbuger issues in VSCode.
    \item Check of Raúl's class notes.
    \item Continue reading Draine's paper on radiative corrections in polarizability \cite{Draine1988ApJ...333..848D}.
    \item Continue with the development of the algorithm to obtain the polarizability of the particles.
    \item Implement the calculation method for Optical Forces given a polarizability, a field, and a field gradient.
    \item Documentation of quasiestatic polarizability theory.
\end{itemize}

\printbibliography

\end{document}
